\documentclass[aspectratio=169,11pt]{beamer}

% ============================================
% PACKAGES
% ============================================
\usepackage[utf8]{inputenc}
\usepackage[T1]{fontenc}
\usepackage[ngerman]{babel}
\usepackage{lmodern}
\usepackage{tcolorbox}
\tcbuselibrary{skins, hooks}
\usepackage{listings}
\usepackage{xcolor}
\usepackage{fontawesome5}

% Modern Font
\IfFileExists{FiraSans.sty}{\usepackage[sfdefault]{FiraSans}}{}
\renewcommand{\familydefault}{\sfdefault}

% ============================================
% COLORS & DESIGN SYSTEM
% ============================================
\definecolor{BgColor}{HTML}{FAFAFA}
\definecolor{Primary}{HTML}{2B3A42}
\definecolor{Accent}{HTML}{E84A5F}
\definecolor{Secondary}{HTML}{3F5765}
\definecolor{CodeBg}{HTML}{F0F0F0}
\definecolor{Success}{HTML}{28A745}

\setbeamercolor{background canvas}{bg=BgColor}
\setbeamercolor{title}{fg=Primary}
\setbeamercolor{frametitle}{bg=Primary, fg=white}
\setbeamercolor{structure}{fg=Accent}
\setbeamercolor{normal text}{fg=Primary}

\setbeamertemplate{navigation symbols}{}
\setbeamertemplate{footline}{%
  \leavevmode%
  \hbox{%
  \begin{beamercolorbox}[wd=.333333\paperwidth,ht=2.25ex,dp=1ex,center]{author in head/foot}%
    \usebeamerfont{author in head/foot}\tiny Falko Himmel
  \end{beamercolorbox}%
  \begin{beamercolorbox}[wd=.333333\paperwidth,ht=2.25ex,dp=1ex,center]{title in head/foot}%
    \usebeamerfont{title in head/foot}\tiny Grundlagen der Praktischen Informatik
  \end{beamercolorbox}%
  \begin{beamercolorbox}[wd=.333333\paperwidth,ht=2.25ex,dp=1ex,right]{date in head/foot}%
    \usebeamerfont{date in head/foot}\insertframenumber{} / \inserttotalframenumber\hspace*{2ex} 
  \end{beamercolorbox}}%
  \vskip0pt%
}

% ============================================
% CUSTOM BOXES
% ============================================
\tcbset{
    modernbox/.style={
        enhanced,
        boxrule=0pt,
        frame hidden,
        colback=white,
        coltitle=white,
        colbacktitle=Secondary,
        fonttitle=\bfseries,
        drop lifted shadow,
        arc=4pt,
        toptitle=2mm,
        bottomtitle=2mm,
        left=2mm,
        right=2mm
    },
    alertbox/.style={
        modernbox,
        colbacktitle=Accent
    }
}

% ============================================
% CODE LISTINGS
% ============================================
\lstdefinelanguage{Haskell}{
  keywords={class, data, deriving, do, else, if, in, import, let, module, of, then, type, where, case},
  keywordstyle=\color{Accent}\bfseries,
  identifierstyle=\color{Primary},
  sensitive=true,
  comment=[l]{--},
  commentstyle=\color{gray}\ttfamily\itshape,
  stringstyle=\color{teal}\ttfamily,
  morestring=[b]"
}

\lstset{
    language=Haskell,
    backgroundcolor=\color{CodeBg},
    basicstyle=\ttfamily\scriptsize,
    breaklines=true,
    frame=single,
    rulecolor=\color{CodeBg},
    frameround=tttt,
    numbers=left,
    numberstyle=\tiny\color{gray},
    xleftmargin=15pt,
    tabsize=4,
    showstringspaces=false
}

% ============================================
% TITLE PAGE
% ============================================
\title{\textbf{Tutorium 4}}
\subtitle{Verzweigungen, Lokale Definitionen \& Tupel}
\institute{Grundlagen der Praktischen Informatik}
\author{}
\date{\today}

\begin{document}

\begin{frame}[plain]
    \titlepage
\end{frame}

% ============================================
% THEORIE
% ============================================
\section{Verzweigungen}

\begin{frame}[fragile]{Verzweigungen in Haskell}
    Haskell bietet verschiedene Wege, um logische Pfade zu steuern.
    
    \begin{columns}[T]
        \begin{column}{0.31\textwidth}
            \begin{tcolorbox}[modernbox, title=\texttt{if-then-else}]
\begin{lstlisting}
signum If x =
  if x < 0
    then "neg"
    else "pos"
\end{lstlisting}
            \end{tcolorbox}
        \end{column}
        
        \begin{column}{0.33\textwidth}
            \begin{tcolorbox}[modernbox, title=Guards (\texttt{|})]
\begin{lstlisting}
signum G x
  | x < 0     = "neg"
  | otherwise = "pos"
\end{lstlisting}
            \end{tcolorbox}
        \end{column}

        \begin{column}{0.33\textwidth}
            \begin{tcolorbox}[modernbox, title=\texttt{case-of}]
\begin{lstlisting}
signum C x = 
  case compare x 0 of
    LT -> "neg"
    _  -> "pos"
\end{lstlisting}
            \end{tcolorbox}
        \end{column}
    \end{columns}
\end{frame}

\section{Lokale Definitionen}

\begin{frame}[fragile]{Scope begrenzen: \texttt{let} \& \texttt{where}}
    Globale Hilfsmethoden ``verschmutzen'' oft den Namensraum. Um Helper-Methoden privat zu halten, nutzen wir lokale Definitionen.
    
    \begin{columns}[T]
        \begin{column}{0.48\textwidth}
            \begin{tcolorbox}[modernbox, title=Mit \texttt{let ... in}]
\begin{lstlisting}
foo x = 
  let helper y = y + 1 
  in helper x
\end{lstlisting}
            \end{tcolorbox}
            \small Die Definition steht \textbf{vor} der Nutzung.
        \end{column}
        \begin{column}{0.48\textwidth}
            \begin{tcolorbox}[modernbox, title=Mit \texttt{where}]
\begin{lstlisting}
foo x = helper x
  where helper y = y + 1
\end{lstlisting}
            \end{tcolorbox}
            \small Die Definition steht \textbf{nach} der Nutzung. (Oft in Haskell bevorzugt).
        \end{column}
    \end{columns}
\end{frame}

\section{Type Synonyms \& Tupel}

\begin{frame}[fragile]{Tupel und Vektoren}
    Mit \texttt{type} kann man Datentypen (z.\,B. Tupel) benennen (Typ-Synonyme). Wir können direkt auf den Tupel-Inhalten Pattern Matching anwenden.
    
    \begin{tcolorbox}[modernbox, title=Beispiel: Kreuzprodukt]
\begin{lstlisting}
type Vector3D = (Double, Double, Double)

crossProduct :: Vector3D -> Vector3D -> Vector3D
crossProduct (x1, y1, z1) (x2, y2, z2) = 
    let cx = y1 * z2 - z1 * y2
        cy = z1 * x2 - x1 * z2
        cz = x1 * y2 - y1 * x2
    in (cx, cy, cz)
\end{lstlisting}
    \end{tcolorbox}
\end{frame}

% ============================================
% LIVE DEMO
% ============================================
\section{Praxis}

\begin{frame}[plain]
    \begin{center}
        \Huge \textbf{Live Demo} \\
        \vspace{0.5cm}
        \large \faLaptopCode~ \texttt{verzweigungen.hs}, \texttt{lokaleDefinitionen.hs}, \texttt{vector.hs} \\
        \vspace{0.3cm}
        \normalsize Wir programmieren Guards, \texttt{where}-Blöcke und Vektor-Operationen zusammen im Code!
    \end{center}
\end{frame}

\end{document}
